\section{Calibration}
\label{sec:calibration}

The coordinates make every trial slice a positive-density law.  Calibration
therefore has a narrower task: select, among those laws, one that explains the
quotes without inventing unsupported detail between or beyond them.  Four
inputs govern that selection: the quoted price bands, basis order, smoothness
penalty, and tail policy.

\subsection{Use price residuals, measured in volatility units}

Suppose the retained quotes are $(k_i,\sigma_i)$ and
$w_i=\sigma_i^2\tau$.  The basic one-slice objective is
\begin{equation}
\sum_i\left(
\frac{c(k_i;\theta)-\Black(k_i,w_i)}
{\partial_\sigma\Black(k_i,w_i)+\eta}
\right)^2
+\omega_{\rm reg}\sum_{n\ge4}n^{2\nu}a_n^2.
\label{eq:objective}
\end{equation}
Here $\partial_\sigma\Black=\phin(d_+)\sqrt\tau$ is normalized Black vega,
the first-order change in price for a one-unit change in volatility, and
$\eta>0$ is a vega floor.  The numerator remains a price difference, so the
optimizer never inverts an implied volatility.  Dividing by vega merely
expresses a central price error in the volatility units used by the market.
In a remote wing, the floor prevents a nearly worthless option from receiving
an unbounded weight.

The ridge starts at $n=4$.  The modes $a_2$ and $a_3$ carry the main
curvature and asymmetry; higher modes describe progressively finer rank
structure.  The factor $n^{2\nu}$ makes fine oscillation increasingly
expensive while leaving the leading smile shape largely data-driven.

An option quote is usually a bid--ask interval.  If
$[p_i^{\rm lo},p_i^{\rm hi}]$ is the normalized price band and $\phi_i$ the
vega normalizer, a band fit uses
\begin{equation}
\frac{(c_i-p_i^{\rm hi})^++(p_i^{\rm lo}-c_i)^+}{\phi_i},
\qquad
\sqrt{0.05}\frac{c_i-p_i^{\rm mid}}{\phi_i}.
\label{eq:band}
\end{equation}
The first residual is zero inside the spread.  The weak second term selects a
stable center when many laws lie inside a wide band.

A volatility haircut $h$ shrinks unreliable band edges toward the mid:
\begin{equation}
\sigma_i^{\rm lo}=\min(\sigma_i^{\rm bid}+h,\sigma_i^{\rm mid}),
\qquad
\sigma_i^{\rm hi}=\max(\sigma_i^{\rm mid},\sigma_i^{\rm ask}-h).
\label{eq:haircut}
\end{equation}
If the original spread is narrower than $2h$, the interval collapses to the
mid.  In the frozen sample this occurs for every featured SPY quote, while
many wider NVDA bands remain active.

\subsection{How many modes the strip can support}
\label{sec:orderguard}

More modes always improve the best in-sample fit; they do not necessarily
improve the inferred density.  At order $N$ the model has $N+1$ body and
endpoint parameters.  A practical effective order is
\begin{equation}
N_{\rm eff}=\max\left(
\min\left(N,\left\lfloor\frac{n_q}{2}\right\rfloor-1\right),
\min(N,6)\right),
\label{eq:orderguard}
\end{equation}
where $n_q$ is the quote count.  Away from the floor, this keeps roughly two
quotes per shape degree of freedom.  The rule controls identification; it is
not a no-arbitrage condition.  A larger order can change digitals and narrow
butterflies between quotes while leaving vanilla residuals almost unchanged.

The quoted ranks make the same point directly.  After a preliminary fit,
record
\[
[u_{\min},u_{\max}]
=\left[\min_iF_X(k_i),\max_iF_X(k_i)\right].
\]
Modes acting mainly outside this interval are chosen by the ridge and the
tail policy, not by quoted options.  Reporting this interval states which part
of the fitted distribution is directly informed by the market.

\subsection{Optimization coordinates}
\label{sec:charts}

Use the endpoint chart
$(\log\lambda_-,\log\lambda_+,a_2,\ldots,a_N)$ with the modes
\eqref{eq:endpointbasis}.  Body moves then leave both tail coefficients fixed.

The right endpoint chart depends on the selected tail exponent:
\begin{equation}
\lambda_+=
\begin{cases}
\expit\xi,&\alpha_+=0,\\
e^\xi,&\alpha_+>0.
\end{cases}
\label{eq:rightchart}
\end{equation}
The first branch makes the finite-forward wall unreachable.  The second is
unconstrained because every positive moment exists.  A separate moment prior
may still cap $\lambda_+$ for numerical or risk purposes.

The raw coefficients $L$ and $R$ follow from the endpoint equations
\eqref{eq:endpoints}.  This conversion is linear once the body coefficients
are fixed.  Analytic derivatives in the endpoint chart are obtained by
multiplying the raw-coordinate Jacobian by this small chart Jacobian.

\subsection{Choose the tail class outside the optimizer}
\label{sec:alphapolicy}

The exponent $\alpha$ changes the asymptotic class, not merely the local wing
slope.  Over any finite strike range, a sufficiently small positive $\alpha$
can resemble $\alpha=0$ while implying that every moment is finite and the
Lee slope is zero.  Listed vanilla options usually identify tail scale much
more strongly than tail class.  Unconstrained per-slice estimation is
therefore ill-conditioned.

The implementation uses one of three policies:
\begin{enumerate}
\item $\alpha_-=\alpha_+=0$ for exponential log-return tails and power gross-
return tails;
\item fixed exponents chosen for an asset class or product, with separate
left and right values allowed;
\item a small grid of exponent scenarios, each recalibrated, with downstream
prices reported across the scenarios.
\end{enumerate}
The preferred practice is to calibrate a small number of stated tail
scenarios and compare the downstream quantities that depend on them.  Within
one underlier, exponents are held common across expiry unless a monotone
term-structure rule is imposed.  This also stabilizes full-line calendar
order; \cref{sec:calendar} gives the exact restriction.

\subsection{Starting values and solve sequence}

For $\alpha_\pm=0$, the constant-speed seed
\begin{equation}
s=\frac{\sqrt{3w_0}}{\pi}
\label{eq:coldstart}
\end{equation}
matches log-return variance.  It is about eight percent low in ATM implied
variance at small $w_0$, so it is only a starting value.  For other tail
exponents, set all body modes to zero and solve one scalar equation for the
common endpoint level that matches the ATM call.  This costs one-dimensional
root finding and avoids an arbitrary variance proxy.

A complete slice fit follows six steps:
\begin{enumerate}
\item infer $F$ and $D$, convert strikes to $k_i$, and convert bid, mid, and
ask volatilities to normalized prices once;
\item fix $(\alpha_-,\alpha_+)$, select $N_{\rm eff}$, and construct the
residual rows and weights;
\item initialize the endpoint chart from the scalar reference slice or from
the preceding expiry;
\item at each step, build $(x,x',G,G')$ on one grid, price every quote with
\eqref{eq:call}, and obtain the Jacobian from
\eqref{eq:pass1}--\eqref{eq:pass3};
\item rebuild the accepted parameters on the publication grid and compute ATM
handles, moments, wing asymptotics, variance-swap strike, and numerical
certificates from that object;
\item for several expiries, replace independent sequential fits by the joint
calendar-constrained solve of \cref{sec:globalcalendarfit}.
\end{enumerate}

Market conventions are fixed before optimization.  Every trial slice remains
a law.  All reported quantities are read from the same final accepted object;
there is no second smoothing or extrapolation layer that can alter its
density, moments, or tails.
